\chapter{结论与展望}
\section{结论}

本文针对完全弱形式下的时空有限元中，时间末端边界条件难以施加的问题，提出时间末端虚位移本质边界条件施加方案，并以此建立完全弱形式下的拉格朗日乘子型时空有限元法。
该方法基于哈密顿原理的伽辽金弱形式，将弱形式中的虚位移作为拉格朗日乘子构建投影约束，并将该约束回代到原有能量泛函中，以此建立全新拉格朗日乘子型的伽辽金弱形式。
不同于其他时空有限元法，该方法在时间和空间维度均为弱形式，并且验证了其原有的欧拉--拉格朗日方程的等价性。
方法的刚度矩阵具有对称性，同时可根据矩阵的复用性进行程序优化。
所提方案在动力问题（时间维度）、“1+1”维（空间1维+时间维度）动力问题及“2+1”维（空间2维+时间维度）波动问题中进行了理论推导与数值验证。
本文的主要研究结论总结如下：

首先，针对时间维度下动力问题，完全弱形式中时间末端虚位移必需为零导致本质边界条件难以数值施加的问题，
提出了一种基于引入独立拉格朗日乘子的时域末端本质边界条件施加方法，
从而构建了全新增广能量泛函的拉格朗日乘子型完全伽辽金弱形式时空混合有限元法，
并证明了该方法的弱形式和原有欧拉--拉格朗日动力学强形式控制方程之间的等价性。
接着用无阻尼弹簧小车和双自由度楔形--滑块动力学系统的算例，
验证了基于拉格朗日乘子变分求解方案的高精度和稳定性，
并且对比了基于微分控制方程时间强形式差分离散的传统Newmark-$\beta$迭代法，
从结果上看，本文所提方案有不仅能计算出精度最为逼近精确解的位移数值、
还能在随时间增长的计算中不产生数值耗散或数值发散的优点。

其次，针对空间1维加上时间维度的“1+1”维的一维动力问题，
提出了将基于拉格朗日乘子的变分求解方案推广至时空全域网格进行联合离散，
通过引入与虚位移等效的拉格朗日乘子构造出全新的双变量能量泛函，
从而建立了完全弱形式的时空混合伽辽金有限元分析的方法。
利用分部积分法证明了该方法与经典的一维动力问题欧拉--拉格朗日偏微分强形式方程的构等价性，
并确定了拉格朗日乘子变量的求解空间，引入罚函数法来施加本质边界条件。
接着采用一些经典的一维杆问题算例验证了该方案的数值精确性与理论误差的最优收敛率，
对比了时空有限元间断伽辽金法，从结果上看，本文所提方案不会出现随着时间步的增加而出现数值振荡现象、
能够在时空多层级并行计算框架下准确计算。
针对任意分布的网格节点计算出现不稳定的现象，
采用了对梯度由大变小的部分做局部加密的非结构化网格；
和网格最长边二分法（BiRefine）局部细化方案，
实现波动方程强梯度区域的高精确追踪。

最后，针对包含二维平面物理空间联合时间维度的“2+1”维平面波动问题，
同样采用了基于拉格朗日变分方案的完全弱形式时空混合有限元法，验证了其与多维波动问题的控制方程的等价性，
再用典型的三维波动问题算例验证了所提方案的数值精确性和计算稳定性。






\section{展望}

尽管本文在完全弱形式的时空混合有限元方法上取得了阶段性的理论进展与模型验证，
但面对未来更加复杂多变的实际工程需求与深层物理演化现象，
本研究仍有进一步拓展和迭代的空间。未来的研究工作可重点围绕以下几个方面展开：

第一，深化完全非结构化与非均布时空网格体系下的算法稳定性研究。
尽管本文通过局部网格重构与加密策略取得了一定成效，
但在全域节点任意分布的复杂非均布网格时，所提的时空弱形式方案在计算表现上仍未能达到绝对稳定。
特别是在应力和位移等物理量发生变化的高梯度区域，
以及受到拉格朗日乘子投影约束的时间末端边界处，
数值解仍会出现数值振荡与不稳定的现象。
未来亟需探寻网格混合离散或其他有效办法，以解决因网格畸变所引发的数值色散现象，
做到完全稳定的强鲁棒计算。


第二，向“3+1”维真实三维实体空间拓展。
本文目前重点推演与验证了二维空间加一维时间的波动控制问题。
未来亟需将该理论模型向完整的三维实体空间(即构建出真正的四维时空域)进行高维映射与全面推广。
随着计算维度的提升，系统节点自由度与时空全局矩阵规模将呈现暴涨趋势，
针对该联合体系下超大型高维稀疏泛矩阵的高性能快速求解器与储存降维压缩策略将是下一步工程化应用的核心挑战。

